% Root research draft. Integration and whole-proof audit remain pending.
\section{Gaussian fourfolds: polarization, dominance, and integral degrees}

Throughout this section $\pi:C\to X$ is a connected etale cyclic cover of
degree four, $g(X)=3$, and $\sigma$ generates its deck group. Write
$J=J(C)$, $\Theta=\Theta_C$, $D=C/\langle\sigma^2\rangle$, and
\[
 B=\operatorname{im}(1-\sigma^2)\subset J,
 \qquad P=\ker(\operatorname{Nm}_{C/X})^0.
\]
Thus $g(C)=9$, $g(D)=5$, $\dim B=4$, and $\dim P=6$.
Intersection degrees in $B$ will be measured against its principal Prym
polarization $\xi$, whereas degrees in $J$ use $\Theta$ or the explicitly
specified very ample class $3\Theta$.

\begin{proposition}[HD-G01: the kernel is the full two-torsion]
\label{HD-G01}
The restricted Jacobian polarization on $B$ equals $2\xi$. Its type is
$(2,2,2,2)$ and its polarization kernel is $B[2]$, of order $2^8$.
\end{proposition}
\begin{proof}
The map $C\to D$ is an etale double cover and its connected norm kernel
is precisely $B$. Apply Mumford, \emph{Prym Varieties I}, Section~3,
Corollary~2, immediately after Corollary~1 on printed p.~332.
That result identifies the restriction, not just its degree: its kernel
is the full two-torsion and it is twice a principal polarization.
The dimensions above then give the asserted type and order.
This is a cited classical result applied to the intermediate cover, not
a new proof of the Prym principal-polarization theorem.
\end{proof}

\begin{proposition}[HD-G06: the two principal Gaussian lattice types]
\label{HD-G06}
The principal Prym component in HD-G04 has even Gaussian Hermitian
lattice $U\oplus U$. The balanced CM point $E^2\times E^2$ with product
principal polarization has the odd lattice $I_{2,2}$. For every
principally polarized Gaussian Weil fourfold with integral maximal-order
action, either lattice type admits the uniform bound
\[
 \delta_W\le D_0(\mathbf Q(i))
       =4D_{\mathrm{Pr}}=3\,656\,994\,324\,480.
\]
For the Prym type the sharper bound of HD-G04 applies. In HD-G05 one
may use either value of $D$ according to this stated normalization.
\end{proposition}
\begin{proof}
We give the lattice details because a principal polarization does not
by itself specify which integral Gaussian component is meant.
For an etale four-cover choose a symplectic marking of the base so
that its character is nontrivial on only one handle. For each of the
other two handles, the anti-invariant homology has the basis
\[
 a_0-a_2,\ a_1-a_3,\ b_0-b_2,\ b_1-b_3,
\]
where subscripts indicate sheets. Its Jacobian intersection form is
twice the standard symplectic form; after passing to $\xi$ it is
the standard form. The deck action sends the first difference to the
second and the second to minus the first, and does the same for the
$b$'s. In the Gaussian basis consisting of the first $a$- and
$b$-differences the Hermitian matrix is
$\bigl(\begin{smallmatrix}0&i\\-i&0\end{smallmatrix}\bigr)$, which is
isometric to $U$ by a unit change of basis. These independent handle
blocks exhaust rank eight over $\mathbf Z$. This proves the first
claim. The product CM polarization instead gives diagonal signs
$(1,1,-1,-1)$ under the balanced action.

For completeness, the elementary classification required here is as
follows. A principal polarization with integral $\mathbf Z[i]$ action
gives a free unimodular Hermitian lattice of signature $(2,2)$.
Its rational Hermitian space is split: this follows from the
signature-and-determinant classification in Schoen (1998), Section~7,
since its determinant is one. Choose a rational isotropic vector,
scale it to a primitive lattice vector $e$, and use unimodularity to
choose $f$ with $H(e,f)=1$. Their span has matrix
$\bigl(\begin{smallmatrix}0&1\\1&k\end{smallmatrix}\bigr)$,
$k\in\mathbf Z$, and is a unimodular orthogonal summand. Replacing
$f$ by $f+ce$, $c\in\mathbf Z[i]$, changes $k$ by an even integer.
An even $k$ gives $U$; an odd $k$ gives $I_{1,1}$.
The complement is again unimodular and rationally split of signature
$(1,1)$, so the same argument applies. Mixed sums of $U$ and an odd
plane are isometric to $I_{2,2}$: on
$\langle1\rangle\oplus U$, with basis $a,e,f$, the integral
unimodular basis $a+e,a-f,a+e-f$ has norms $1,1,-1$ and is orthogonal.
Thus the only possibilities are the even $U^2$ and the odd $I_{2,2}$.

Finally, with $S=\bigl(\begin{smallmatrix}1&1\\1&-1\end{smallmatrix}\bigr)$,
\[
 S^*I_{1,1}S=2U,\qquad S^*US=2I_{1,1}.
\]
On two blocks this yields a fixed Gaussian-linear isogeny from an
even-type fourfold to any prescribed odd-type point, after transporting
the complex structure by the rational similitude. Its degree is
$2^4=16$ and its pullback of the target principal polarization is
$2\xi$. A detecting surface pushes forward to a detecting integral
cycle, and one reduced component detects. Its degree is at most
$2^2D_{\mathrm{Pr}}$. This proves the all-principal bound, with no
unbounded isogeny transport and no appeal to the core index $J$.
\end{proof}

\begin{proposition}[HD-G02: published rank four of the Gaussian Prym period map]
\label{HD-G02}
On a nonempty open subset of the moduli of connected etale cyclic
four-covers of genus-three curves, the differential of the period map
to the four-dimensional Gaussian Weil family has rank exactly four.
Consequently it dominates its polarized integral-lattice component.
\end{proposition}
\begin{proof}
The exact published assertion is van Geemen (1996),
\emph{Theta functions and cycles on some abelian fourfolds},
Theorem~5.3(3), printed pp.~629--630. It proves this Gaussian variant
of Schoen's argument. The following calculation records the source
check and its integral-component interpretation; the rank-four
dominance theorem itself is not claimed as a new result of this run.
Let $\eta$ be the associated line bundle of exact order four on $X$.
The cotangent map is the multiplication map
\[
 \mu:H^0(X,K_X\eta)\otimes H^0(X,K_X\eta^{-1})
       \longrightarrow H^0(X,K_X^2).
\]
The formula is the cyclic-character restriction of Schoen (1988),
Lemma~3.11, printed pp.~29--30; see also Koike (2004), equation~(1),
printed p.~568, for the same order-four character pair.
Both factors on the left have dimension two by Riemann--Roch and
$\eta^{\pm1}\ne\mathcal O_X$.

We verify the order-four variant of the argument on Schoen's printed
p.~30; his Theorem~3.2 itself is stated for the Eisenstein field.
If $K_X\eta$ has a base point $p$, Riemann--Roch gives
$h^0(\eta^{-1}(p))=1$, so $\eta^{-1}=\mathcal O_X(q-p)$ for some
$q\ne p$. Hence $4q\sim4p$. A rational function with divisor
$4q-4p$ defines a degree-four map $X\to\mathbf P^1$ totally ramified
at $p$ and $q$. Riemann--Hurwitz gives total ramification twelve;
these two points account for six. There are therefore at most six
other branch values. The corresponding Hurwitz loci, after the two
distinguished values are placed at zero and infinity and the remaining
one-dimensional automorphism group is removed, have dimension at
most five. There are finitely many ramification types and Nielsen
classes in degree four with this ramification bound. Their images
cannot contain an open subset of the six-dimensional $\mathcal M_3$.
Thus $K_X\eta$ is base-point-free for a general pair $(X,\eta)$.

The base-point-free pencil trick now identifies $\ker\mu$ with
$H^0(X,\eta^{-2})=0$, because $\eta$ has exact order four.
Thus $\operatorname{rank}\mu=2\cdot2=4$, and the differential has the
same rank. The target component has dimension four, giving dominance.
This supplies the actual Gaussian rank calculation; the source locator
for Schoen's method and the distinct field in his printed theorem are
kept separate.
\end{proof}

\begin{proposition}[HD-G03: the norm fibre needed in the Bezout argument]
\label{HD-G03}
The fibre of $\operatorname{Nm}_{C/X}:\operatorname{Pic}^4(C)
\to\operatorname{Pic}^4(X)$ over $K_X$ is a union of translates of the
sixfold $P$, not translates of the fourfold $B$. Each connected
component $T_0$ has
\[
 \operatorname{type}(\Theta|_P)=(1,1,1,1,4,4),\qquad
 \deg_{3\Theta}T_0=3^6\,6!\,16=8\,398\,080.
\]
The full fibre $W_4(C)\cap\operatorname{Nm}^{-1}(K_X)$ has dimension
at most two.
\end{proposition}
\begin{proof}
The norm is surjective, and its connected kernel has dimension
$g(C)-g(X)=6$. The polarization type is the etale cyclic-cover
formula of Lange--Ortega, \emph{Prym varieties of cyclic coverings},
Section~4, first paragraph on the induced polarization (preprint
arXiv:0805.1020, p.~6), citing Birkenhake--Lange, Corollary~12.1.4 and
Lemma~12.3.1. In their notation the type is
$(1^{(d-2)(g-1)},d^{g-1})$; here $d=4,g=3$.
Riemann--Roch for polarized abelian varieties gives the displayed
degree. The morphism $\operatorname{Sym}^4 C\to\operatorname{Sym}^4 X$
is finite. The inverse image of $|K_X|\simeq\mathbf P^2$ therefore
has dimension two. Its image under the Abel map is exactly the
indicated intersection, which proves the dimension assertion.

In particular the four-dimensional degree
$3^4\,4!\,16$ cannot be substituted for the degree of this norm fibre.
The polarization calculation in HD-G01 does not justify that
substitution: it concerns a different abelian subvariety.
\end{proof}

\begin{lemma}[HD-G04a: an effective supplier component detects after projection]
\label{HD-G04a}
Some irreducible surface component $Y$ of
$W_4(C)\cap T_0$, for one connected norm-fibre component $T_0$,
has the property that $(1-\sigma^2)_*[Y]$ has nonzero Weil projection
in $B$.
\end{lemma}
\begin{proof}
Let $s:C^4\to J$ be the sum of Abel maps, and put
$u=(1-\sigma^2)s:C^4\to B$.
For each primitive character $\chi$ of the cyclic group, the
determinant line in the tensor product of four copies of
$H^1(C)_\chi$ is the pullback of the corresponding Weil line on $B$.
The alternating-tensor map is injective. Products of two classes of
the same primitive character in $H^1(C)$ vanish, since their product
would have nontrivial character $\chi^2$ in the trivial representation
$H^2(C)$. It follows that $u^*W$ lies in the all-$H^1$ Kunneth
summand of $H^4(C^4)$ and is primitive for the product polarization.
It is a real Hodge subspace of type $(2,2)$; the middle pairing on it
is positive definite by Hodge--Riemann.

Schoen (1988), Theorem~2.0, its unramified construction on printed
pp.~12--13, and Lemma~2.6 show that this determinant subspace is
spanned by the classes of signed combinations of surfaces supported
over $|K_X|$ in $C^4\to\operatorname{Sym}^4 X$. No ramified-cover
repair is used here. For nonzero $w\in W$, the nonzero class $u^*w$
therefore pairs nontrivially with at least one of the effective
irreducible supplier surfaces $Z$. By the projection formula
$u_*[Z]$ pairs nontrivially with $w$.

Put $Y=s(Z)$. This image must have dimension two: otherwise $u_*[Z]$
vanishes in the relevant dimension. The pushforward $s_*[Z]$ is a
positive integer times $[Y]$, so $(1-\sigma^2)_*[Y]$ still has
nonzero pairing. The support lies in the norm fibre because
$Z$ lies over $|K_X|$. By HD-G03 that fibre intersected with $W_4(C)$
has dimension at most two; hence $Y$ is an irreducible component.
The Weil projector is orthogonal for the middle pairing on this
primitive summand, so the nonzero pairing is precisely detection.
\end{proof}

\begin{theorem}[HD-G04: an explicit detector bound on the principal Prym component]
\label{HD-G04}
For the Gaussian principal Prym component $S_{\mathrm{Pr}}$, with
polarization $\xi$ and integral $\mathbf Z[i]$ action,
\[
 \delta_W\le D_{\mathrm{Pr}}=914\,248\,581\,120
\]
on every fibre.
\end{theorem}
\begin{proof}
The Poincare formula gives
\[
 \deg_{3\Theta}W_4(C)=3^4\frac{9!}{5!}=244\,944.
\]
For the detecting component $Y$ supplied by HD-G04a, the projective
Bezout bound for components of an improper intersection (Fulton,
\emph{Intersection Theory}, second edition, Example~8.4.6) gives
\[
 \deg_{3\Theta}Y\le244\,944\cdot8\,398\,080.
\]
Set $p=1-\sigma^2:J\to B$. The involution $\sigma^2$ preserves
$\Theta$ and is self-adjoint for its Rosati involution. Thus
$p^\dagger p=2(1-\sigma^2)$ has eigenvalues zero and four, so
$p^*(\Theta|_B)\le4\Theta$ as semipositive invariant forms.
Since $\Theta|_B=2\xi$, this says $p^*\xi\le2\Theta$.
The image $p(Y)$ is a surface and is still detecting; its reduced
image has no greater degree than its effective pushforward. Therefore
\[
 \deg_\xi p(Y)
 \le \int_Y(p^*\xi)^2
 \le4\deg_\Theta Y
 =\frac49\deg_{3\Theta}Y
 \le914\,248\,581\,120.
\]
All cycles here are reduced integral subvarieties or integral proper
pushforwards; no division of an effective cycle was performed.
HD-G02 supplies a dense set of fibres of the component. The supplied
detector closedness result (division attack, Theorem~D, equivalently
Lemma~Dom together with the Chow argument) extends the same integer
bound to the whole component.
\end{proof}

\section{A quadratic refinement of the division step}

\begin{theorem}[HD-WQ: fixed quadratic multiplication removes a Gram factor]
\label{HD-WQ}
Use the setting and projector data of the supplied Theorem~W, with
$R$ replaced by its absolute value and $\kappa$ its stated integral
degree mass. Put
\[
 q=\operatorname{Nm}(a)^n,\quad
 s=a^{2n}+\bar a^{2n}\in\mathbf Z,\quad H=4q^2-s^2>0.
\]
For a detector of integral degree $D$, a nonzero rational Weil class
$t$, and $\delta t$ integral, let $T=(-1)^n\int t^2>0$.
There exist integral effective positive and negative cycles representing
$Nt$, with
\[
 N\le\delta H\frac{R^2}{(n!)^2}D^2,
 \qquad
 \deg Z^\pm\le
 \frac{4\delta q\sqrt{HT}\,\kappa R}{n!}D^2.
\]
Ceilings make these into an explicit integral bound. In particular
the supplied quartic comparison may be sharpened to a quadratic
comparison; this does not by itself bound $D$ in a family.
\end{theorem}
\begin{proof}
Use, without re-proving, Steps~1--4 of Theorem~W. They give an
integral $y=R e_W[Y]\ne0$ represented with effective positive and
negative parts of degrees at most $\kappa D$, and
\[
 d=\lVert y\rVert^2=(-1)^n\int y^2\in\mathbf Z_{>0},
 \qquad \sqrt d\le RD/n!.
\]
For $A=[a]^*|_W$, put $y_1=y$ and $y_2=Ay$.
The second generator has parts of degrees at most $q\kappa D$.
The Rosati adjoint and the two eigenvalues give
\[
 A+A^\dagger=s\operatorname{id},\qquad
 A^\dagger A=q^2\operatorname{id}.
\]
Consequently the Gram matrix is exactly
\[
 G=d\begin{pmatrix}1&s/2\\s/2&q^2\end{pmatrix}.
\]
This fixed shape is the point that permits cancelling a common
factor from the original Cramer calculation.
Let $b_j=\langle\delta t,y_j\rangle\in\mathbf Z$, with the positive
middle pairing on $W$. Direct inversion gives
\[
 \delta Hd\,t=(4q^2b_1-2sb_2)y_1+(4b_2-2sb_1)y_2.
\]
Its multiplier $N=\delta Hd$ is a positive integer even when $s$ is odd.
The two coefficient functionals are pairing with
$4q^2y_1-2sy_2$ and $4y_2-2sy_1$. Their squared norms, computed
from the displayed Gram matrix, are respectively $4q^2Hd$ and
$4Hd$. Cauchy--Schwarz bounds their absolute values by
$2\delta q\sqrt{HTd}$ and $2\delta\sqrt{HTd}$.
Forming the indicated integral linear combination of the signed
cycle representatives gives each effective part degree at most
\[
 \kappa D\bigl(2\delta q\sqrt{HTd}
             +q\,2\delta\sqrt{HTd}\bigr)
 \le\frac{4\delta q\sqrt{HT}\,\kappa R}{n!}D^2.
\]
The bound for $N$ follows from the same estimate on $d$.
There is no cancellation by a rational coefficient in the cycle group;
the common factor was removed from the integer Gram formula before
the cycles were assembled.
\end{proof}

\begin{theorem}[HD-G05: an integral Gaussian fourfold division constant]
\label{HD-G05}
Let $(A,\xi,\mathbf Z[i])$ be a principally polarized Gaussian Weil
fourfold possessing a detector of degree at most the integer $D$.
For $\delta t$ integral and $T=\int t^2>0$, put $h=\lceil\sqrt T\rceil$.
Then
\[
 c_t(A)\le C_4(D;\delta,T)
 =\delta D^2\max\{36\,864,76\,800h\}.
\]
In particular $D=D_{\mathrm{Pr}}$ gives a fully explicit constant on
the component in HD-G04.
\end{theorem}
\begin{proof}
Write $I=[i]^*$ and $S=[1+i]^*$ on $H^4(A)$.
On the character pieces indexed by $j=0,1,2,3,4$, the eigenvalues
of $I$ are $(-1)^j$ and those of $S$ are
$-4,-4i,4,4i,-4$, respectively (the order of conjugate entries is
immaterial). Hence
\[
 (1+I)(4-S)=16e_W.
\]
For a detector $Y$ of degree at most $D$, the integral class
\[
 y=(4+4I-S-IS)[Y]
\]
has effective positive and negative parts of degrees at most $8D$:
$I$ preserves degree and $S$ multiplies surface degree by four.
The orthogonal-projector and comass estimates already proved in
Theorem~W give $d=\int y^2\le64D^2$.

Let $A_*= [2+i]^*|_W$. Its eigenvalues are $-7\pm24i$, and
$y_2=A_*y$ has positive and negative degrees at most $200D$.
The Gram matrix of $(y,y_2)$ is
\[
 d\begin{pmatrix}1&-7\\-7&625\end{pmatrix}.
\]
For $b_j=\langle\delta t,y_j\rangle\in\mathbf Z$ there is the exact
integral identity
\[
 576\delta d\,t=(625b_1+7b_2)y+(7b_1+b_2)y_2.
\]
The dual vectors $625y+7y_2$ and $7y+y_2$ have norms
$600\sqrt d$ and $24\sqrt d$. Thus $N=576\delta d$ is at most
$36\,864\delta D^2$, while each effective part of the displayed
representative has degree at most
\[
 8D\,(600+25\cdot24)\,\delta\sqrt{Td}
 \le76\,800\delta\sqrt T\,D^2.
\]
Replacing $\sqrt T$ by its integer ceiling proves the statement.
\end{proof}
